%%
%% spring-lecture12.tex
%% 
%% Made by Alex Nelson <pqnelson@gmail.com>
%% Login   <alex@lisp>
%% 
%% Started on  2026-04-30T10:41:23-0700
%% Last update 2026-04-30T10:41:23-0700
%% 

\lecture[Indecomposable modules]{}

\begin{definition}
Let $A$ be an algebra. An $A$-module $M$ is called
\define{Indecomposable} if $M\neq0$ and if we cannot decompose it into
a nontrivial direct sum: $M=P\oplus Q$ then either
$P=0$ or $Q=0$.

Otherwise, if there exists such a $P$ and $Q$ such that $M=P\oplus Q$,
we say that $M$ is \define{Decomposable}.
\end{definition}

(Simple modules are indecomposable; we should view indecomposable
modules as a generalization of simple modules. The main result of this
focus will be the Krull--Schmidt theorem.)

\begin{caution}
By convention, the zero module is neither decomposable nor indecomposable.
\end{caution}

\begin{proposition}\label{prop:5.1}
Let $A$ be an algebra, $M$ be an $A$-module.
If $M$ is Noetherian or Artinian, then $M$ is a finite direct sum of
indecomposable submodules.
\end{proposition}

\begin{proof}
Without loss of generality assume $M\neq0$. If $N$ is a minimal
nonzero direct summand of $M$, then $N$ is indecomposable. We claim
such an $N$ exists [Pf: If $M$ is Artinian, then obviously; if $M$ is
Noetherian, then $N$ is a complement of a maximal proper direct
summand of $M$.]
By repeating this process, $M$ decomposes as
\begin{equation}
M=N_{1}\oplus M_{1}=N_{1}\oplus N_{2}\oplus M_{2}=\dots
\end{equation}
whree $N_{1}$ is the minimal direct summand of $M$ and $M_{1}$ is the
``remainder', $N_{2}$ is the minimal direct summand of $M_{1}$ and
$M_{2}$ is its ``remainder'', and so on. Since
\begin{equation}
M_{1}\supset M_{2}\supset\dots\quad\mbox{(descending)}
\end{equation}
and
\begin{equation}
N_{1}\subset N_{1}\oplus N_{2}\subset\dots\quad\mbox{(ascending)}
\end{equation}
there exists some $k$ such taht $M=N_{1}\oplus\cdots\oplus N_{k}$.
\end{proof}

\begin{definition}
We say the algebra $A$ is a \define{Local Algebra} if $A/\Jacobson(A)$
is a division algebra.
\end{definition}

\begin{node}
If $A$ is a local algebra, then $1_{A}\neq0$ (i.e., $A$ is nontrivial).
\end{node}

\begin{node}
The next proposition (and its corollary) give a sort of analog to
how we characterize simple modules via Schur's lemma.
\end{node}

\begin{proposition}\label{prop:5.2}
For a nonzero algebra $A$, the following are equivalent:
\begin{enumerate}
\item $A$ is a local algebra
\item $A\setminus A^{\times}\subset\Jacobson(A)$
\item $A\setminus A^{\times}$ is closed under addition.
\end{enumerate}
\end{proposition}

\begin{proof}
$(1)\implies(2)$ For any $x\notin\Jacobson(A)$, there exists some
$y\in A$ such that $xy-1\in\Jacobson(A)$ and also $yx-1\in\Jacobson(A)$.
By Proposition~\ref{prop:4.9},
\begin{equation}
xy=1+(xy-1)\in A^{\times},
\end{equation}
and in a similar manner
\begin{equation}
yx=1+(yx-1)\in A^{\times}.
\end{equation}
So $x\in A^{\times}$. [Pf: $xy(xy)^{-1}=1$ and $(yx)^{-1}yx=1$ are the
right and left inverses.] Hence $x\notin\Jacobson(A)$ implies $x\in A^{\times}$.

\bigbreak
$(2)\implies(3)$ For $x\in\Jacobson(A)$, if $x\in A^{\times}$ then by
Proposition~\ref{prop:4.9} $0=1+x(-x)^{-1}\in A^{\times}$ which is a
contradiction because $A\neq0$. So if $x\in\Jacobson(A)$, then
$x\notin A^{\times}$ (i.e., $A\setminus A^{\times}\supset\Jacobson(A)$),
so by hypothesis (2) $A\setminus A^{\times}=\Jacobson(A)$ is closed
under addition since it is a two-sided ideal.

\bigbreak$(3)\implies(1)$
For $x\notin\Jacobson(A)$, by Proposition~\ref{prop:4.9} there exists
$y,z\in A$ such that $1+xy\notin A^{\times}$ and $1+zx\notin A^{\times}$.
This means $xy\in A^{\times}$ and $zx\in A^{\times}$ by hypothesis (3)
[Pf: let $w_{1}=1+xy$ and $w_{2}=1+zx$, then $w_{1}-w_{2}\in A^{\times}$
therefore $xy\in A^{\times}$ and $zx\in A^{\times}$ otherwise $1\in A\setminus A^{\times}$ by (3)].
Thus $x\in A^{\times}$ in the same manner. Therefore
$A\setminus\Jacobson(A)\subset A^{\times}$ and so $A/\Jacobson(A)$ is
a division algebra.
\end{proof}

\begin{corollary}\label{cor:5.3}
\begin{enumerate}
\item If all non-units of $A$ are nilpotent, then $A$ is a local algebra
\item If $N$ is an $A$-module such that $\End_{A}(N)$ is a local
algebra, then $N$ is indecomposable.
\end{enumerate}
\end{corollary}

\begin{proof}
\begin{enumerate}
\item Let $0\neq x\notin A^{\times}$ be a non-unit of $A$. Suppose $k\in\NN$
is the smallest positive number such that $x^{k}=0$. Then $xy\notin A^{\times}$
for any $y\in A$, otherwise $x^{k-1}(xy)=0$ so $x^{k-1}=0$ which
contradicts the minimality of $k$. Then any element in $xA$ are
nilpotent, so $x\in xA\subset\Jacobson(A)$ by Corollary~\ref{cor:4.10}~(2).
This means $A\setminus A^{\times}\subset\Jacobson(A)$ thus $A$ is
local by Proposition~\ref{prop:5.2}.
\item Assume $\End_{A}(N)$ is local. Without loss of generality
$\id_{N}\neq0$ so $N\neq0$. If $N=P\oplus Q$, then we can consider
projections $\pi\colon N\to P$ and $\rho\colon N\to Q$. Then
$\pi+\rho=\id_{N}$ and since $\End_{A}(N)$ is local either $\pi$ or
$\rho$ is in $(\End_{A}N)^{\times}$ by Proposition~\ref{prop:5.2}. But
$\pi^{2}=\pi$ and $\rho^{2}=\rho$ so it follows that either
$\pi=\id_{N}$ or $\rho=\id_{N}$ (since
$\pi\circ\pi^{-1}=\id=(\pi^{2})\circ\pi^{-1}=\pi$ and similarly for
$\rho$). This means either $Q=0$ or $P=0$ (respectively), which
implies $N$ is indecomposable.\qedhere
\end{enumerate}
\end{proof}

\begin{xca}[{Pierce~\cite[\S5.2 Ex~1]{pierce1982associative}}]
Prove an algebra $A$ is local if and only if $A$ has a unique maximal
right ideal.
\end{xca}

\subsection{Fitting's Lemma}

\begin{lemma}
Let $A$ be an algebra, let $M$ be an $A$-module.
For $\phi\in\End_{A}M$, the following all holds:
\begin{enumerate}
\item If $M$ is Noetherian and $\phi$ is surjective, then
  $\phi\in\Aut_{A}(M)$ is an automorphism.
\item If $M$ is Artinian and $\phi$ is injective, then
  $\phi\in\Aut_{A}(M)$ is an automorphism.
\end{enumerate}
\end{lemma}

\begin{proof}
\begin{enumerate}
\item Since $0\subset\ker(\phi)\subset\ker(\phi^{2})\subset\dots$,
  there exists an $n\in\NN$ such that $\ker(\phi^{n})=\ker(\phi^{n+1})=\dots$.
  In particular, the preimage of
\begin{subequations}
  \begin{align}
(\phi^{n})^{-1}(\ker(\phi)) &= (\phi^{n+1})^{-1}(0)\\
&=\ker(\phi^{n+1})\\
    &=\ker(\phi^{n})\\
    &=(\phi^{n})^{-1}(0).
  \end{align}
\end{subequations}
Since $\phi$ is surjective by hypothesis, then so too is
$\phi^{n}$. And so
\begin{subequations}
  \begin{align}
\ker(\phi) &=\phi^{n}(\phi^{n})^{-1}(\ker(\phi))\quad\mbox{by surjectivity}\\
&=\phi^{n}(\phi^{n})^{-1}(0)\\
&=(0)
  \end{align}
\end{subequations}
hence $\phi$ is injective, so $\phi$ is an automorphism.
\item Homework.\qedhere
\end{enumerate}
\end{proof}

\begin{lemma}[Fitting's lemma]\label{lemma:5.5}
Let $A$ be an algebra. Let $M$ be an $A$-module which is both Artinian
and Noetherian. For any $\phi\in\End_{A}M$ there exists a
decomposition of $M=P\oplus Q$ such that
\begin{enumerate}
\item $\phi(P)\subset P$ and
\item $\phi(Q)\subset Q$ and
\item $\phi|_{P}\in\Aut_{A}P$ and $\phi|_{Q}$ is nilpotent.
\end{enumerate}
\end{lemma}

\begin{proof}
We see
$M\supset\phi(M)\supset\phi^{2}(M)\supset\dots$ is a descending chain
and $0\subset\ker(\phi)\subset\ker(\phi^{2})\subset\dots$ is an
ascending chain. The Artinian and Noetherian conditions on $M$ implies
there exists an $m\in\NN$ such that $\phi^{m}(M)=\phi^{n}(M)$ and
$\ker(\phi^{m})=\ker(\phi^{n})$ for all $n\geq m$. Denote
\begin{equation}
P=\phi^{m}(M)\quad\mbox{and}\quad Q=\ker(\phi^{m}).
\end{equation}
then
\begin{equation}
\phi(P)=P,
\end{equation}
and
\begin{subequations}
  \begin{align}
\phi(Q) &=\phi(\ker(\phi^{m}))\\
&=\phi(\ker(\phi^{m+1}))\\
&\subset\ker(\phi^{m})=Q
  \end{align}
\end{subequations}
so $\phi(Q)\subset Q$, and also
\begin{equation}
\phi^{m}(Q)=\phi^{m}(\ker(\phi^{m}))=0,
\end{equation}
which implies $\phi|_{Q}$ is nilpotent. Moreover
\begin{equation}
P\cap Q=0,
\end{equation}
since $\phi|_{P\cap Q}$ is injective and nilpotent.

Finally, we need to check $M=P\oplus Q$. We see that
\begin{subequations}
  \begin{align}
M &= (\phi^{m})^{-1}(\phi^{m}(M))\\
&=(\phi^{m})^{-1}(\phi^{2m}(M))\\
&=(\phi^{m})^{-1}(\phi^{2m}(M)+0)\\
&=\phi^{m}(M)+\ker(\phi^{m})\\
&=P+Q.
  \end{align}
\end{subequations}
Hence the claim.
\end{proof}

\begin{corollary}\label{cor:5.6}
Let $A$ be an algebra. Let $M$ be a Noetherian and Artinian $A$-module.
Then $\End_{A}M$ is local if and only if $M$ is indecomposable.
\end{corollary}

\begin{proof}
\forwardproof\ By Corollary~\ref{cor:5.3}~(2).

\backwardproof\ By Lemma~\ref{lemma:5.5} every element of $\End_{A}M$ is either
nilpotent or a unit. Then by Corollary~\ref{cor:5.3}~(1) $\End_{A}M$
is local.
\end{proof}


